\documentclass[11pt]{article}
\usepackage[margin=1in]{geometry}
\usepackage{amsmath,amssymb,booktabs}
\title{An elementary prime-count balance lemma, a scale ladder with
reproducible numerics, and a gated record of a dielectric sphere's
mass-energy invariants}
\author{Author Name}
\date{2026-09-25}

\begin{document}
\maketitle

\begin{abstract}
We prove an elementary lemma on the balance \(S(n)=(n-1)-2\pi(n)\), establish
an 18-rung length/scale ladder with a gate script reproducing every asserted
value, and give mass/energy bookkeeping of a canonical 550 nm dielectric
sphere. Every numerical statement here is reproduced by one script
(\texttt{results\_of\_record.py}, 51 checks, exit 0). The speculative framing
layer that motivates the ladder is separated in Part II and is explicitly
non-claimant.
\end{abstract}

\section{Part I: Results}

\subsection{The prime-count balance lemma}
Let \(S(n)=(n-1)-2\pi(n)\). Then
\[
S(n)<0 \iff n\in\{2,\dots,8\},\qquad
S(n)=0 \iff n\in\{1,9,11,13\},
\]
\(n=10\) is the first strict composite majority, and \(S(n)\ge 0\) for
every \(n\ge 9$.
\begin{proof}[Proof sketch]
For \(n\le 16\) the assertion is verified directly from
\(\pi(9..16)=4,4,5,5,6,6,6,6\). For \(n\ge 17\), Rosser--Schoenfeld gives
\(\pi(n)<1.25506\,n/\ln n\), hence
\(2\pi(n)<2.51012\,n/\ln n\le n-1\), since \(2.51012/\ln n\le 1-1/n\) at
\(n=17\) and the left side decreases. Thus \(S(n)\ge 0\) for all \(n\ge 17\).
\qed
\end{proof}

\subsection{The 18-rung scale ladder}
For a canonical source diameter \(D=550\) nm and rung radius \(s\),
\(n(u)=D/s\) is asserted for every rung of the table below (check 30 of the
gate).

\begin{center}
\begin{tabular}{lrr}
\toprule
Rung & \(\log_{10}(s/\text{m})\) & \(n(u)=D/s\) \\
\midrule
Planck & $-34.79$ & $3.40\times10^{28}$ \\
quark & $-19.00$ & $5.5\times10^{12}$ \\
proton & $-15.08$ & $6.55\times10^8$ \\
atom & $-10.00$ & $5.5\times10^3$ \\
molecule & $-9.00$ & $5.5\times10^2$ \\
virus & $-7.00$ & $5.5$ \\
ball & $-6.26$ & $1.0$ \\
cell & $-5.00$ & $0.055$ \\
human & $0.24$ & $3.14\times10^{-7}$ \\
neutron star (PSR J0740+6620) & $4.09$ & $4.44\times10^{-11}$ \\
Earth & $6.80$ & $8.62\times10^{-14}$ \\
Sun & $9.14$ & $3.95\times10^{-16}$ \\
Kuiper 50 AU & $12.87$ & $7.35\times10^{-20}$ \\
quasar broad-line region (100-day lag) & $15.41$ & $2.12\times10^{-22}$ \\
galaxy 30 kpc & $20.97$ & $5.94\times10^{-28}$ \\
Great Attractor 50 Mpc & $24.19$ & $3.56\times10^{-31}$ \\
Laniakea 160 Mpc & $24.69$ & $1.11\times10^{-31}$ \\
observable universe (particle horizon) & $26.64$ & $1.25\times10^{-33}$ \\
\bottomrule
\end{tabular}
\end{center}

The log10 gaps of this ladder are all distinct (0.51--15.79, mean 3.614,
sd 3.504, CV 0.970) --- descriptive, not a significance claim. The
neutron-star size is the equatorial radius of PSR J0740+6620
(R$=$12.39\,km, Riley et al.\ 2021); the quasar size is $c\tau$ for a
100-day reverberation lag (Kaspi et al.\ 2000).

\subsection{The canonical sphere}
With \(\rho=1100\,\text{kg/m}^3\), \(R=275\) nm,
\[
M=\rho\,\tfrac43\pi R^3=9.583\times10^{-17}\;\text{kg},\quad
E=Mc^2=8.612\;\text{J},\quad
r_s=\frac{2GM}{c^2}=1.423\times10^{-43}\;\text{m}.
\]
The sphere sits \(\approx36.3\) orders above its own Schwarzschild radius
(check 29); its Bekenstein bound is \(4.71\times10^{20}\,k_B\), about
\(8.2\times10^9\) times the order-of-magnitude thermal entropy (check 34).

\subsection{On the observable ball}
The observable causal 3-ball at the current particle-horizon radius
\(R=4.4\times10^{26}\) m (Hubble-parameter band
\({67.4\le H_0\le 73.04}\) \(\Rightarrow\) \(R\in[4.06,4.4]\times10^{26}\) m)
has volume \(V=\tfrac43\pi R^3\in[2.80,3.568]\times10^{80}\ \text{m}^3\) and
holds \(3.2\times10^{99}\)--\(4.1\times10^{99}\) canonical balls (checks 27, 32).
Two observers whose horizons are separated by \(R/2\) share \(63.28\%\) of
each ball's volume (check 33).

\section{Part II: Ideas (non-claimant)}%
\label{ideas}
This section is the author's vocabulary, not physics. The zero-to-horizon
recursion, the "Law of Object Zeros", and the reading of mass as a
gauge-invariant residue are framing in the corpus's internal `Law Apprehension
of Works` thread; they are carried here only so that reviewers see exactly
what is not claimed. Nothing in Part II alters Part I.

\begin{thebibliography}{9}
\bibitem{rosser} J.~B.~Rosser, L.~Schoenfeld, {Approximate formulas for some
functions of prime numbers}, Illinois J. Math. 6 (1962) 64--94.
\bibitem{planck} Planck Collaboration, A&A 641 (2020) A6.
\bibitem{riess} A.~Riess et al., ApJ 934 (2022) L7; 938 (2022) 36.
\end{thebibliography}

\end{document}